\chapter{Rates and FX}\label{ch:iv2:rates-and-fx}

\begin{ivtrains}
Curves and what you can pull from them, hedging with DV01, swaps and futures, short-rate models and swaptions, then FX forwards and options. Unless a stem says otherwise, the yields are annually compounded, and every level here is invented for the arithmetic.
\end{ivtrains}

\section{Curves at speed}

\begin{question}[$\star$ \iqroles{QR, QT, Trader, Risk} \iqfirm{bank, hedge fund} \iqlevel{junior}]\label{iq:iv2:rates-and-fx:1}
A five-year zero-coupon bond yields 4\%, with annual compounding. What is its price per \$100? What is the price if 4\% is continuously compounded?
\end{question}

\begin{question}[$\star$ \iqroles{QR, QT, Trader, Risk} \iqfirm{bank, hedge fund} \iqlevel{junior}]\label{iq:iv2:rates-and-fx:2}
A one-year bond pays a coupon of \$5 and \$100 at maturity, and its yield is 4\%. What is its price? Without computing, say what its price would be at a yield of 5\%, and why.
\end{question}

\begin{question}[$\star$ \iqroles{QR, QT, Trader, Risk} \iqfirm{bank, hedge fund} \iqlevel{junior}]\label{iq:iv2:rates-and-fx:3}
A five-year zero yields 4\%. What are its Macaulay and modified durations? By how much does a \$100 price change if yields rise 10 basis points?
\end{question}

\begin{question}[$\star$ \iqroles{QR, QT, Trader, Risk} \iqfirm{bank, hedge fund} \iqlevel{junior}]\label{iq:iv2:rates-and-fx:4}
You own \$10 million of a bond trading at par with a modified duration of 8. What is your DV01? How much do you make if yields fall 25 basis points, ignoring convexity?
\end{question}

\begin{question}[$\star$ \iqroles{QR, QT, Trader, Risk} \iqfirm{bank, hedge fund} \iqlevel{junior}]\label{iq:iv2:rates-and-fx:5}
One-year zeros yield 3\% and two-year zeros 4\%. What is the one-year rate one year from now, as a forward? What does a borrower do to lock it in today?
\end{question}

\begin{question}[$\star$ \iqroles{QT, Trader, Risk} \iqfirm{bank, hedge fund} \iqlevel{junior}]\label{iq:iv2:rates-and-fx:6}
You are long \$10 million of a five-year bond with a modified duration of 4.5, at par. You hedge with a ten-year bond with a modified duration of 8, also at par. How much do you short? What are you still exposed to?
\end{question}

\begin{question}[$\star$ \iqroles{QR, QT, Trader, Risk} \iqfirm{bank, hedge fund} \iqlevel{junior}]\label{iq:iv2:rates-and-fx:7}
You pay fixed at 4\% on a \$10 million five-year swap, and receive floating. It has a DV01 of about \$4,500. Do you gain or lose if rates rise? What bond position is this like? What do you make if the swap's rate rises 10 basis points?
\end{question}

\begin{question}[$\star$ \iqroles{QT, Trader, Analyst} \iqfirm{bank, hedge fund} \iqlevel{junior}]\label{iq:iv2:rates-and-fx:8}
A three-month bill trades at a discount of 4\%, on the 90/360 basis. What is its price per \$100? What is the simple money-market yield on the same basis? Why is it above 4\%?
\end{question}

\begin{question}[$\star$ \iqroles{QT, Trader, QR} \iqfirm{bank, hedge fund} \iqlevel{junior}]\label{iq:iv2:rates-and-fx:9}
EURUSD is 1.10. The dollar one-year rate is 5\%, and the euro one-year rate is 3\%, both simple. What is the one-year forward? Quote the forward points. Does the forward predict that the euro rises?
\end{question}

\begin{question}[$\star$ \iqroles{QT, Trader, Analyst} \iqfirm{bank, hedge fund} \iqlevel{junior}]\label{iq:iv2:rates-and-fx:10}
EURUSD is 1.10 and USDJPY is 150. What is EURJPY? If EURJPY is quoted at 166, what do you do? What do you need to check before calling it a free profit?
\end{question}

\begin{question}[$\star$ \iqroles{QT, Trader, QR} \iqfirm{bank, hedge fund} \iqlevel{junior}]\label{iq:iv2:rates-and-fx:11}
You hold a call on 1 million euros against the dollar, at a strike of 1.10. What option is that, seen from the other side, and on how many dollars? At what strike in euros per dollar?
\end{question}

\begin{question}[$\star$ \iqroles{QT, Trader, QR} \iqfirm{bank, hedge fund} \iqlevel{junior}]\label{iq:iv2:rates-and-fx:12}
You borrow euros at 3\% and lend dollars at 5\%, unhedged, with EURUSD at 1.10. Where does EURUSD have to be in a year for you to break even, with simple rates? How is that related to the forward? Why does the trade's name tell you nothing about its risk?
\end{question}

\begin{question}[$\star$ \iqroles{QT, Trader, QR, Risk} \iqfirm{bank, hedge fund} \iqlevel{junior}]\label{iq:iv2:rates-and-fx:13}
A bond has a modified duration of 8 and a convexity of 80. Estimate the percentage price change for a rise and for a fall of 100 basis points. Which is larger, and why do people like that?
\end{question}

\section{Hedges and curve trades}

\begin{question}[$\star\star$ \iqroles{QR, QT, QD} \iqfirm{bank, hedge fund} \iqlevel{junior}]\label{iq:iv2:rates-and-fx:14}
One-year and two-year annual par yields are 3\% and 4\%. Bootstrap the discount factors. What is the two-year zero rate? What is the forward rate between years one and two?
\end{question}

\begin{question}[$\star\star$ \iqroles{QR, QT, QD} \iqfirm{bank, hedge fund} \iqlevel{junior}]\label{iq:iv2:rates-and-fx:15}
Zero rates at one, two and three years are 3\%, 4\% and 5\%. What is the three-year par swap rate, with annual payments? What is the two-year swap starting in a year? Why does the floating leg not need a model?
\end{question}

\begin{question}[$\star\star$ \iqroles{QT, Trader, QR} \iqfirm{bank, hedge fund} \iqlevel{junior, experienced}]\label{iq:iv2:rates-and-fx:16}
On a flat 4\% curve, you compare a five-year 4\% par bond with a barbell of a two-year and a ten-year 4\% par bond, with the same duration. How are the two weighted? What are their convexities? What does each make on a parallel move of 100 basis points up or down? So why does anyone hold the bullet, given that a real curve is not flat?
\end{question}

\begin{question}[$\star\star$ \iqroles{QT, Trader} \iqfirm{bank, hedge fund} \iqlevel{junior, experienced}]\label{iq:iv2:rates-and-fx:17}
Three par bonds of two, five and ten years have a DV01, per \$100, of 1.886, 4.452 and 8.111 cents. You short \$100 of the five-year and buy the two and ten so that the trade has no cash and no DV01. What is the size of each? What is your profit if the five-year yield is up 10 basis points and the other two do not move? What are you betting on?
\end{question}

\begin{question}[$\star\star$ \iqroles{QT, Trader, Risk} \iqfirm{bank, hedge fund} \iqlevel{junior}]\label{iq:iv2:rates-and-fx:18}
The two-year has a DV01 of \$190 per \$1 million and the ten-year \$800 per \$1 million. Set up a DV01-neutral steepener with \$10 million of the ten-year. What do you make if the two-year falls 5 basis points and the ten-year rises 5? What if both rise 10?
\end{question}

\begin{question}[$\star\star$ \iqroles{QT, Trader, QR, Risk} \iqfirm{bank, hedge fund} \iqlevel{junior, experienced}]\label{iq:iv2:rates-and-fx:19}
Yield changes have an annualised volatility of 80 basis points for the two-year and 75 for the ten-year, with a correlation of 0.7. You hold \$10 million of the ten-year, with a DV01 of \$8,000, and you hedge with the two-year. What DV01 of hedge does a regression give? How does the residual risk compare with a DV01-neutral hedge?
\end{question}

\begin{question}[$\star\star$ \iqroles{QR, QT, Risk} \iqfirm{bank, hedge fund} \iqlevel{junior}]\label{iq:iv2:rates-and-fx:20}
You receive fixed at 4\% on a swap with two annual payments left, on a notional of \$100. The curve is now flat at 3\%. What is the swap worth? Get it two ways: as a bond less par, and as the rate difference times the annuity.
\end{question}

\begin{question}[$\star\star$ \iqroles{QR, QT, Trader} \iqfirm{bank, hedge fund} \iqlevel{junior, experienced}]\label{iq:iv2:rates-and-fx:21}
A three-month interest rate future expires in five years, and it fixes on a rate from year 5 to year 5.25. The rate's normal volatility is 1\% a year. What is the convexity adjustment between the future's rate and the forward rate, using the usual rule? Which is higher, and why?
\end{question}

\begin{question}[$\star\star$ \iqroles{QT, Trader, QR} \iqfirm{bank, hedge fund} \iqlevel{junior, experienced}]\label{iq:iv2:rates-and-fx:22}
A bond future trades at 100. The short can deliver one of three bonds, with conversion factors of 0.92, 1.10 and 1.25, and prices of 93.0, 110.8 and 125.6. Which is the cheapest to deliver? Why does the answer change when yields move? Ignore carry.
\end{question}

\begin{question}[$\star\star$ \iqroles{QT, Trader, QR} \iqfirm{bank, hedge fund} \iqlevel{junior, experienced}]\label{iq:iv2:rates-and-fx:23}
You buy a five-year par bond at a yield of 4\%, financed at a repo rate of 3\%. The four-year yield is 3.8\%, and a four-year bond has a duration of 3.6. What is your carry and your roll-down over a year if the curve does not change? How far can the four-year yield rise at the horizon before you lose money?
\end{question}

\begin{question}[$\star\star$ \iqroles{QR, QD, Risk} \iqfirm{bank, hedge fund} \iqlevel{junior, experienced}]\label{iq:iv2:rates-and-fx:24}
In Vasicek's model, the short rate starts at 3\%, reverts to 5\% with a speed of 0.5, and has a volatility of 1\%. What are the mean and the standard deviation of the short rate in two years? What is the chance of a negative rate then? Does that make the model useless?
\end{question}

\begin{question}[$\star\star$ \iqroles{QR, QD} \iqfirm{bank, hedge fund} \iqlevel{junior, experienced}]\label{iq:iv2:rates-and-fx:25}
A short rate has no drift, so it stays flat in expectation at a level $f$ (Ho-Lee with zero drift), and it has a 1\% normal volatility. What is the ten-year zero yield in basis points below $f$? Derive the $\sigma^2T^2/6$. Why does a volatile rate lower long yields?
\end{question}

\section{Options on rates, and FX}

\begin{question}[$\star\star$ \iqroles{QR, QT, Trader} \iqfirm{bank, hedge fund} \iqlevel{junior}]\label{iq:iv2:rates-and-fx:26}
Price a caplet on a three-month rate fixing in a year, with a forward rate of 4\%, a strike of 4\%, a lognormal volatility of 20\% and a discount factor of 0.95 to the payment date, per unit of notional. What is the same volatility in basis points, and the caplet by Bachelier? What is the normal volatility if the rate is at 2\%, and at the same lognormal volatility?
\end{question}

\begin{question}[$\star\star$ \iqroles{QR, QT, Trader} \iqfirm{bank, hedge fund} \iqlevel{junior}]\label{iq:iv2:rates-and-fx:27}
An at-the-money one-year swaption is on a swap with an annuity of 4.5. The normal volatility is 80 basis points. What is it worth as a percentage of notional, by the rule of thumb? Why does the annuity multiply it? What does one basis point of volatility change the price by?
\end{question}

\begin{question}[$\star\star$ \iqroles{QR, QT, Trader} \iqfirm{bank, hedge fund} \iqlevel{junior}]\label{iq:iv2:rates-and-fx:28}
A one-year payer swaption and a one-year receiver swaption have the same strike of 5\%. The forward swap rate is 4\%, and the annuity is 4.5. What is the payer less the receiver? Prove it. What is the position, and which model did you use?
\end{question}

\begin{question}[$\star\star$ \iqroles{QT, Trader, QR} \iqfirm{bank, hedge fund} \iqlevel{junior, experienced}]\label{iq:iv2:rates-and-fx:29}
EURUSD spot is 1.10, the one-year euro rate is 3\%, and the one-year forward trades at 1.1245. What dollar rate does the swap market imply, on simple rates? The dollar cash rate is 5\%. What does this difference tell you? Why does it not get arbitraged away?
\end{question}

\begin{question}[$\star\star$ \iqroles{QR, QT, Trader} \iqfirm{bank, hedge fund} \iqlevel{junior}]\label{iq:iv2:rates-and-fx:30}
EURUSD is 1.10. A one-year at-the-money-spot euro call has a strike of 1.10, a volatility of 10\%, a dollar rate of 5\% and a euro rate of 3\%, continuously compounded. What is the forward? What are the call and the put worth, in dollars per euro? What are the spot delta and the forward delta?
\end{question}

\begin{question}[$\star\star$ \iqroles{QR, QT, Trader} \iqfirm{bank, hedge fund} \iqlevel{junior, experienced}]\label{iq:iv2:rates-and-fx:31}
A one-year euro call, with EURUSD at 1.10, a strike of 1.10, a 10\% volatility, a dollar rate of 5\% and a euro rate of 3\%, continuously compounded, is worth \$0.0536 per euro. What is that in euros, and as a percentage of the euro notional and of the dollar notional? Show that the same option priced from the dollar side, as a dollar put on the rate in euros per dollar at a strike of 0.909, gives the same value. What must the volatility be, and why?
\end{question}

\begin{question}[$\star\star$ \iqroles{QT, Trader, QR} \iqfirm{bank, hedge fund} \iqlevel{junior, experienced}]\label{iq:iv2:rates-and-fx:32}
An FX desk quotes the one-year smile as at-the-money 10\%, a 25-delta risk reversal of $-2\%$ and a 25-delta butterfly of 0.5\%. What are the 25-delta call and put volatilities? With a forward of 1.1222, a euro rate of 3\%, and using the spot delta, what strikes are they at?
\end{question}

\begin{question}[$\star\star$ \iqroles{QT, Trader, QR, Risk} \iqfirm{bank, hedge fund} \iqlevel{junior, experienced}]\label{iq:iv2:rates-and-fx:33}
EURUSD has a volatility of 10\%, and USDJPY a volatility of 9\%, with a correlation of 0.5 between their log changes. What is the volatility of EURJPY? What if the correlation is $-0.5$? How do you use this in reverse, and what does it mean for trading three currencies?
\end{question}

\section{Models, products and conventions}

\begin{question}[$\star\star\star$ \iqroles{QR, QT, Trader} \iqfirm{bank, hedge fund} \iqlevel{junior, experienced}]\label{iq:iv2:rates-and-fx:34}
An overnight swap's floating leg compounds a daily rate over a week: 4\%, 4\%, 5\%, 4\% on Monday to Thursday, and 4\% on Friday for three days, on an actual/360 basis. What rate does the leg pay for the week? How does that compare with a simple average? Unlike a term rate set in advance, when is it known? What does that do to a caplet on it?
\end{question}

\begin{question}[$\star\star\star$ \iqroles{QR, QD} \iqfirm{bank, hedge fund} \iqlevel{junior, experienced}]\label{iq:iv2:rates-and-fx:35}
Write Vasicek's zero-coupon price in closed form, with a speed of 0.5, a mean of 5\%, a volatility of 1\% and a start of 3\%. What is the five-year zero and its yield? Derive $B(T)$. Check the price with a Monte Carlo simulation of the discounting. How does the yield behave at very long maturity?
\end{question}

\begin{question}[$\star\star\star$ \iqroles{QR, QD} \iqfirm{bank, hedge fund} \iqlevel{experienced}]\label{iq:iv2:rates-and-fx:36}
In Hull and White's one-factor model, with a speed of 0.1 and a volatility of 1\%, the market curve is flat at 4\%, continuously compounded. Find the drift $\theta(t)$ that fits it. Why is Vasicek's constant mean not enough? Check by simulation that the model gives back the ten-year zero. What is $f(t,T)$ later on?
\end{question}

\begin{question}[$\star\star\star$ \iqroles{QR, QT, Trader} \iqfirm{bank, hedge fund} \iqlevel{experienced}]\label{iq:iv2:rates-and-fx:37}
Two annual forward rates start at 4\%, with lognormal volatilities of 20\%, and they fix in a year, with discount factors of 0.96 and 0.92. Take the swap rate as their average. A cap has two caplets with a 4\% strike, and a payer swaption has the same strike. What is each worth in basis points, with correlations of 1, 0.5 and 0? Why is the swaption never worth more than the cap? What market does this link?
\end{question}

\begin{question}[$\star\star\star$ \iqroles{QR, QT, Risk} \iqfirm{bank, hedge fund} \iqlevel{junior, experienced}]\label{iq:iv2:rates-and-fx:38}
A five-year 5\% bond can be called at par by the issuer in one year. Take a flat yield, a normal yield volatility of 100 basis points, and value the call as an option on the four-year bond a year out. What are the bullet, the call and the callable worth at yields of 4\%, 5\% and 6\%? What are their durations? What is negative convexity, and who is short it?
\end{question}

\begin{question}[$\star\star\star$ \iqroles{QT, Trader, QR} \iqfirm{bank, hedge fund} \iqlevel{experienced}]\label{iq:iv2:rates-and-fx:39}
A future can deliver a five-year or a twenty-year 6\% bond, each with a conversion factor of 1, since the factors are computed at a 6\% yield. The yield is at 6\%, a flat curve, and at delivery it is normal with a standard deviation of 100 basis points. What is the future's fair price, as the mean of the cheaper of the two? How much does the short's choice take off? Which bond is cheapest if yields rise?
\end{question}

\begin{question}[$\star\star\star$ \iqroles{QR, QT, Risk} \iqfirm{bank, hedge fund} \iqlevel{junior, experienced}]\label{iq:iv2:rates-and-fx:40}
Yield changes at 2, 5, 10 and 30 years have volatilities of 80, 85, 75 and 65 basis points. Their correlations are 0.85, 0.7 and 0.6 for the 2-year with the others, 0.9 and 0.8 for the 5-year with the longer two, and 0.9 for the 10-year with the 30-year. What share of variance do the first three principal components explain? What do they look like? How would you use them to hedge?
\end{question}

\begin{question}[$\star\star\star$ \iqroles{QR, QT} \iqfirm{bank, hedge fund} \iqlevel{experienced}]\label{iq:iv2:rates-and-fx:41}
A contract pays an annual rate that fixes in five years, and it pays at the fixing and not a year later. The forward is 4\%, with a lognormal volatility of 20\%. What is the convexity adjustment? Derive it from a change of measure, and compare the exact value with the first-order rule. Why does it arise?
\end{question}

\begin{question}[$\star\star\star$ \iqroles{QR, QT} \iqfirm{bank, hedge fund} \iqlevel{experienced}]\label{iq:iv2:rates-and-fx:42}
A constant-maturity swap pays the ten-year swap rate, fixed in five years, in a single payment at the fixing. The forward swap rate is 4\%, with a lognormal volatility of 20\%. Use a flat-yield annuity $A(y)$ to estimate the adjustment. What does the first-order rule give? What does a simulation under the annuity measure give? Which way does it go, and why does it depend on the whole smile?
\end{question}

\begin{question}[$\star\star\star$ \iqroles{QR, QT} \iqfirm{bank, hedge fund} \iqlevel{experienced}]\label{iq:iv2:rates-and-fx:43}
A dollar-paid contract is on a yen index at 100, paying one dollar per index point. The index has a volatility of 20\%, the dollar value of a yen has a volatility of 10\%, and their correlation is $-0.3$. Rates are zero. What is the quanto forward? Derive the drift adjustment, and check it by simulation. Why does the sign go that way?
\end{question}

\begin{question}[$\star\star\star$ \iqroles{QT, Trader, QR} \iqfirm{bank, hedge fund} \iqlevel{experienced}]\label{iq:iv2:rates-and-fx:44}
A one-year FX option on a forward of 1.1222 has a 10\% volatility. What is the strike at which a straddle has zero delta, when the delta does not include the premium? When it does? Why do some currency pairs use the premium-adjusted form? What goes wrong if you confuse the two?
\end{question}

\begin{question}[$\star\star\star$ \iqroles{QD, QR, Risk} \iqfirm{bank, hedge fund} \iqlevel{junior, experienced}]\label{iq:iv2:rates-and-fx:45}
Continuously compounded zero rates at 1, 2, 3 and 5 years are 3\%, 4\%, 3.8\% and 4.2\%. You interpolate linearly in the zero rate. What is the instantaneous forward either side of two years? And before and after three years? What does log-linear interpolation of discount factors give? Which would you hedge with, and why?
\end{question}

\begin{question}[$\star\star\star$ \iqroles{QR, QT, Trader} \iqfirm{bank, hedge fund} \iqlevel{junior, experienced}]\label{iq:iv2:rates-and-fx:46}
A five-year par bond has a duration of 4.45. With a normal yield volatility of 80 basis points, what is the bond's price volatility? Value a one-year at-the-money-forward call on it by the rule of thumb. Compare this with a one-year receiver swaption at 80 basis points on an annuity of 4.5. Why are they nearly the same? Which of the two carries the extra risk?
\end{question}

\begin{question}[$\star\star\star$ \iqroles{QT, Trader, QR} \iqfirm{bank, hedge fund} \iqlevel{junior, experienced}]\label{iq:iv2:rates-and-fx:47}
A bond priced at 100 pays a coupon of 2 in the next half-year, and it is the cheapest to deliver into a future, with a future's price times the conversion factor of 99.5 at delivery in half a year. What is the implied repo rate, on a simple basis? The market repo for that bond is 2.5\%. What do you trade, and what can spoil it?
\end{question}

\section{Three to build}

\begin{question}[$\star\star\star\star$ \iqroles{QR, QD} \iqfirm{bank, hedge fund} \iqlevel{experienced}]\label{iq:iv2:rates-and-fx:48}
In the Hull-White model with a speed of 0.1, a volatility of 1\% and a flat 4\% continuously compounded curve, price an at-the-money one-year payer swaption on a five-year annual swap. Use Jamshidian's decomposition, and check by simulation. What is the strike? Why does the decomposition work in one factor, and fail in two?
\end{question}

\begin{question}[$\star\star\star\star$ \iqroles{QR, QT, QD} \iqfirm{bank, hedge fund} \iqlevel{experienced}]\label{iq:iv2:rates-and-fx:49}
A one-year FX smile has quotes at-the-money 10\%, at a strike of 1.1278, 25-delta put 11.5\% at 1.0482, and 25-delta call 9.5\% at 1.1992. The forward is 1.1222, and the dollar rate is 5\%. Derive the vanna-volga weights that replicate a new option with those three. What does it give at strikes of 1.10 and 1.25? Why do the pillars come back exactly? Where should you stop trusting it?
\end{question}

\begin{question}[$\star\star\star\star$ \iqroles{QR, QD} \iqfirm{bank, hedge fund} \iqlevel{experienced}]\label{iq:iv2:rates-and-fx:50}
In Hull and White's model with a speed of 0.1, a volatility of 1\% and a flat 4\% continuously compounded curve, a payer Bermudan swaption can be exercised at years 1, 2, 3 and 4, into a swap with annual payments that ends at year 5, at a strike equal to the fair rate of the swap starting at year 1. Roll it back on a grid. What are the four European values and the Bermudan? What is the switch value? What does the speed of mean reversion do to that switch value, and why?
\end{question}
